\begin{frame}[fragile]
\frametitle{Motivace}

\begin{columns}
\begin{column}{0.45\textwidth}
Algoritmus A\\
\begin{minted}{c}
int matrix[N][N];
int main() {
  long int i, j, sum1 = 0;
  for(i=0; i<N; i++)
    for(j=0; j<N; j++)
      sum1 += matrix[i][j];
}
\end{minted}
\end{column}
\hfill
\begin{column}{0.45\textwidth}
Algoritmus B\\
\begin{minted}{c}
int matrix[N][N];
int main() {
  long int i, j, sum1 = 0;
  for(i=0; i<N; i++)
    for(j=0; j<N; j++)
      sum1 += matrix[j][i];
}
\end{minted}
\end{column}
\end{columns}
\bigskip
Oba dva programy vypadají velmi podobně. 

Program A prochází pole po řádcích, program B prochází pole po sloupcích.
\bigskip

Bude se nějak lišit doba výpočtu?

\end{frame}

\begin{frame}[fragile]
\frametitle{Motivace}

\begin{columns}
\begin{column}{0.45\textwidth}
Algoritmus A\\
\begin{minted}{c}
int matrix[N][N];
int main() {
  long int i, j, sum1 = 0;
  for(i=0; i<N; i++)
    for(j=0; j<N; j++)
      sum1 += matrix[i][j];
}
\end{minted}
\end{column}
\hfill
\begin{column}{0.45\textwidth}
Algoritmus B\\
\begin{minted}{c}
int matrix[N][N];
int main() {
  long int i, j, sum1 = 0;
  for(i=0; i<N; i++)
    for(j=0; j<N; j++)
      sum1 += matrix[j][i];
}
\end{minted}
\end{column}
\end{columns}
\bigskip

\begin{tabular}{|l|l|l|}\hline
N & A & B \\\hline
100000 & 12.791328s &138.047563s \\\hline
10000 & 0.126945s &0.486535s \\\hline
1000 & 0.001329s &0.001756s \\\hline
100 & 0.000083s &0.000094s \\\hline
\end{tabular}
\end{frame}

\begin{frame}
\frametitle{Co je paměť}

\centering

\includegraphics[width=.95\linewidth]{generated/memory/addr-space-cz.pdf}

\end{frame}

\begin{frame}
\frametitle{Co je paměť -- popis}

Paměť :
\begin{itemize}
\item je pole adresovatelných buněk
\begin{itemize}
\item buňky mají shodnou šíři - většinou 8 bitů (byte) nebo jejich násobek
\end{itemize}
\item velikost fyzického adresového prostoru je omezena počtem bitů adresy, kterou je CPU schopno nastavit na paměťové sběrnici, tedy kolik bitů lze použít k adresování paměti. Pro organizaci po bytech pak
\begin{itemize}
\item 16 bitů adresy umožňuje maximálně 64KiB velkou paměť
\item 32 bitů adresy umožňuje maximálně 4GiB velkou paměť
\item 37 bitů adresy (maximum procesoru Intel Core  i9-13900K, pak záleží ještě na základní desce) maximálně 128 GiB velkou paměť
\end{itemize}
\item typ přístupu:
\begin{itemize}
\item k vnitřní paměti počítače lze přistupvat v náhodném pořadí (= random access -- RAM)
\item některé externí paměti (např. magnetopásková paměť využívaná pro zálohování -- levnější než HDD a SSD) umožňují pouze sekvenční přístup, tedy čtení a zápis bajtů po sobě
\item pro HDD, SSD, Flash je možný přítup po celých blocích, u SSD/Flash je zápis bloku možný vždy jednou do smazání po velkých celcích
\end{itemize}
\end{itemize}

\end{frame}

\begin{frame}
\frametitle{Polovodičové paměti -- terminologie}

\begin{description}
 \item[Adresa] vstup, který vybírá svojí hodnotou jednu konkrétní buňku paměti
 \item[Hodnota] vlastní informace
 \item[Stavová data] volitelně další informace (třeba o platnosti hodnoty, redundanci pro opravení chyb, apod.)
\end{description}

Parametry pamětí:

\begin{description}
  \item[Vybavovací doba paměti] (latence) kritický parametr. Délka časového intervalu mezi objevením se požadavku a okamžikem, kdy jsou data k dispozici.
  \item[Doba přístupu] vybavovací doba + obnovení obsahu po destruktivním čtení připadně doba, kdy lze zadat další požadavek
  \item[Propustnost] výkonový parametr. Schopnost zpracovat uvedené množství za jednotku času.
\end{description}

\end{frame}
